\clearpage
\thispagestyle{empty}

\begin{center}
	\vspace*{4cm}
	{\Huge \textbf{KẾT THÚC PHẦN 2}} \\[1.5cm]
	
	\rule{0.8\textwidth}{0.4pt} \\[0.8cm]
	
	\parbox{0.9\textwidth}{
		\centering
		\textit{
			Phần 2 bàn về xử lý dữ liệu, huấn luyện và triển khai các mô hình ngôn ngữ
			lớn: cách dùng các thư viện, và cách tổ chức quy trình MLOps quanh chúng.
			Cùng với lý thuyết ở Phần 1, đó là phần kiến thức cần để xây dựng các ứng
			dụng NLP thật.
		}
	} \\[1cm]
	
	\rule{0.8\textwidth}{0.4pt} \\[2cm]
\end{center}

\noindent
\textbf{Tóm lược nội dung Phần 2:}
\begin{itemize}[leftmargin=*]
	\item \textbf{Chương 1 -- Quy trình làm việc và công cụ xử lý dữ liệu:}
	Phác thảo vòng đời dự án NLP tinh gọn, các kỹ thuật thu thập dữ liệu (APIs, Web Scraping), 
	làm sạch và gán nhãn (Pandas, Polars, Label Studio, Snorkel, LLM-based), 
	tăng cường dữ liệu, và quản lý phiên bản với DVC.
	
	\item \textbf{Chương 2 -- Huấn luyện và đánh giá mô hình:}
	Hệ sinh thái Hugging Face (\texttt{transformers}, \texttt{datasets}, \texttt{accelerate}, \texttt{evaluate}),
	xây dựng các hàm đánh giá thực tế (\texttt{compute\_metrics}), theo dõi thí nghiệm
	(Weights \& Biases, MLflow), huấn luyện phân tán với DeepSpeed,
	toàn bộ chuỗi hậu huấn luyện với TRL (SFT/QLoRA, DPO, GRPO),
	và đánh giá tái lập được với LM Evaluation Harness.
	
	\item \textbf{Chương 3 -- Tối ưu hóa, triển khai và vận hành (MLOps):}
	Cơ sở dữ liệu vector, các thư viện tối ưu hóa (\texttt{bitsandbytes}, \texttt{peft}, AWQ/GPTQ, ONNX),
	đóng gói và phục vụ mô hình (FastAPI, Docker), các framework chuyên dụng (vLLM với prefix caching
	và giải mã suy đoán, SGLang, BentoML, Triton), và vòng đời MLOps hoàn chỉnh (giám sát, tái huấn luyện).
	
	\item \textbf{Chương 4 -- Công thức xây dựng các ứng dụng phổ biến (Recipes):}
	Hướng dẫn từng bước để xây dựng các ứng dụng thực tế: tìm kiếm ngữ nghĩa,
	hệ thống RAG, fine-tuning cho phân loại và NER, xây dựng tác tử đơn giản,
	và trích xuất thông tin có cấu trúc.
\end{itemize}